\chapter{Cadre Conceptuel et Théorie de Risque de Crédit}

\newpage

L'étude de l'évaluation du risque de crédit au sein d'une institution bancaire nécessite, au préalable, une maîtrise rigoureuse des concepts théoriques et réglementaires qui la sous-tendent. Ce premier chapitre a pour objectif de jeter les bases académiques indispensables à la compréhension de notre problématique centrale. Avant de confronter ces modèles à la réalité opérationnelle du Centre d'Affaires, il convient de définir les paramètres fondamentaux de la prise de risque et les outils de modélisation existants dans la littérature financière. Pour ce faire, ce chapitre s'articulera autour de trois axes principaux :

Dans une première section, nous aborderons les fondements théoriques du risque de crédit. Il s'agira de définir la nature de ce risque dans le cadre de l'intermédiation financière et d'analyser l'impact de l'asymétrie d'information (sélection adverse et aléa moral) sur la décision d'octroi de crédit.

La deuxième section sera consacrée au cadre prudentiel international. Une attention particulière sera portée aux évolutions des Accords de Bâle et à l'approche de notation interne (IRB - Internal Ratings-Based), qui dictent les normes de calcul des exigences en fonds propres et institutionnalisent le recours au rating.

Enfin, la troisième section détaillera l'architecture mathématique et statistique des principaux modèles de prévision de défaillance. Nous examinerons les modèles econométriques classiques et statistiques traditionnels, les modèles de reduced form ( Machine Learning/ Models Stochastique), ainsi que les modèles structurels, établissant ainsi le référentiel théorique qui servira de base à notre analyse comparative ultérieure.


\section{La théorie du risque de crédit}
